% Part: lambda-calculus
% Chapter: lambda-definability
% Section: primitive-recursive-functions

\documentclass[../../../include/open-logic-section]{subfiles}

\begin{document}

\olfileid[pt]{lam}{ldf}{prf}
\olsection{As funções recursivas primitivas são \usetoken{S}{lambda definable}} 

Lembre-se de que as funções recursivas primitivas são aquelas que podem ser
definidas a partir das funções básicas $\Zero$, $\Succ$ e $\Proj{n}{i}$
por composição e recursão primitiva.

\begin{lem}
  \ollabel{lem:basic}
  As funções recursivas primitivas básicas $\Zero$, $\Succ$ e as
  projeções~$\Proj{n}{i}$ são !!{lambda definable}.
\end{lem}

\begin{proof}
Elas são !!{lambda define}d pelos seguintes termos:
\begin{align*}
  \fn{Zero} & \ident \lambd[a][\lambd[fx][x]]\\
  \fn{Succ} & \ident \lambd[a][\lambd[fx][f (a f x)]] \\
  \fn{Proj}^n_i & \ident \lambd[x_0\dots x_{n-1}][x_i]
\end{align*}
\end{proof}

\begin{lem}
  \ollabel{lem:comp} Suponha que a função $k$-ária~$f$ e as funções $n$-árias
  $g_0$, \dots,~$g_{k-1}$ sejam !!{lambda definable} pelos termos
  $F$, $G_0$, \dots, $G_{k-1}$, e que $h$ seja definida a partir delas por
  composição. Então $h$ é !!{lambda definable}.
\end{lem}

\begin{proof}
  $h$ pode ser !!{lambda define}d pelo termo
  \[
  H \ident \lambd[x_0 \dots x_{n-1}][F\, (G_0 x_0 \dots
    x_{n-1}) \dots (G_{k-1} x_0 \dots x_{n-1})]
  \]
  Deixamos a verificação desse fato como exercício.
\end{proof}

\begin{prob}
  Complete a prova de \olref[lam][ldf][prf]{lem:comp} mostrando
  que $H\num{n_0}\dots\num{n_{n-1}} \red \num{h(n_0, \dots,
    n_{n-1})}$.
\end{prob}

Observe que \olref{lem:comp} não exigiu que $f$ e $g_0$,
\dots,~$g_{k-1}$ fossem recursivas primitivas; exige-se apenas que
sejam totais e !!{lambda definable}.

\begin{lem}\ollabel{lem:prim}
  Suponha que $f$ seja uma função $n$-ária e~$g$, uma função $n+2$-ária,
  que elas sejam !!{lambda definable} pelos termos $F$ e~$G$, e que a
  função~$h$ seja definida a partir de $f$ e $g$ por recursão
  primitiva. Então $h$ também é !!{lambda definable}.
\end{lem}

\begin{proof}
  Lembre-se de que $h$ é definida por
  \begin{align*}
    h(x_1, \dots, x_n, 0) &= f(x_1, \dots, x_n)\\
    h(x_1, \dots, x_n, y+1) & = g(x_1, \dots, x_n, y, h(x_1, \dots, x_n, y)).
  \end{align*}
  Informalmente, a definição recursiva primitiva itera a
  aplicação da função $g$ $y$ vezes e a aplica a $f(x_1,
  \dots, x_n)$. Isso lembra a definição dos numerais de Church,
  que também são definidos como um iterador.

  Para simplificar, damos a definição e a prova para um único
  argumento adicional~$x$. A função $h$ é !!{lambda define}d por:
  \begin{align*}
    H \ident & \lambd[x][\lambd[y][\fn{Snd} (y
        D \tuple{\num{0}, F x})]]
    \intertext{onde }
    D \ident & \lambd[p][\tuple{\fn{Succ} (\fn{Fst}\, p), (G x
          (\fn{Fst}\, p) (\fn{Snd}\, p))}]
  \end{align*}
  O estado da iteração que mantemos é um par, cujo primeiro componente é o
  $y$ atual e cujo segundo componente é o valor correspondente de~$h$. A
  cada passo da iteração, criamos um par de novos valores de $y$ e
  $h$; ao fim da iteração, devolvemos a segunda parte do
  par, que é o valor final de $h$. Mostraremos agora que esta é de fato uma
  representação da recursão primitiva.

  Queremos provar que, para quaisquer $n$ e $m$, $H\,\num{n}\,\num{m} \red
  \num{h(n,m)}$. Para isso, mostramos primeiro que, se $D_n \ident
  \Subst{D}{\num{n}}{x}$, então $D_n^m \tuple{\num{0}, F\, \num{n}} \red
  \tuple{\num{m}, \num{h(n, m)}}$. Procedemos por indução em~$m$.

  Se $m=0$, queremos $D_n^0 \tuple{\num{0}, F\, \num{n}} \red
  \tuple{\num{0}, \num{h(n, 0)}}$. Mas $D_n^0 \tuple{\num{0}, F\,
    \num{n}}$ é simplesmente $\tuple{\num{0}, F\, \num{n}}$. Como $F$
  !!{lambda define}s~$f$, isso se reduz a $\tuple{\num{0},
    \num{f(n)}}$ e, como $f(n) = h(n, 0)$, isso é $\tuple{\num{0},
    \num{h(n,0)}}$.

  Agora suponha que
  \[
    D_n^m \tuple{\num{0}, F\, \num{n}} \red
    \tuple{\num{m}, \num{h(n, m)}}.
  \]
  Queremos mostrar que
  \[
    D_n^{m+1} \tuple{\num{0}, F\, \num{n}} \red
    \tuple{\num{m+1}, \num{h(n, m+1)}}.
  \]
  \begin{align*}
    D_n^{m+1} \tuple{\num{0}, F\, \num{n}}
    & \ident D_n(D_n^{m} \tuple{\num{0}, F\, \num{n}})\\
    & \red D_n\,\tuple{\num{m}, \num{h(n, m)}} \text{\qquad (pela HI)}\\
    & \ident (\lambd[p][
      \tuple{
        \fn{Succ} (\fn{Fst}\, p), (G \, \num{n} (\fn{Fst}\, p) (\fn{Snd}\, p))
    }]) \tuple{\num{m}, \num{h(n,m)}}\\
    & \redone
    \tuple{\fn{Succ} (\fn{Fst}\, \tuple{\num{m}, \num{h(n,m)}}),\\
      & \qquad (G \, \num{n} (\fn{Fst}\, \tuple{\num{m}, \num{h(n,m)}}) (\fn{Snd}\, \tuple{\num{m}, \num{h(n,m)}}))}\\
    & \red 
    \tuple{\fn{Succ}\, \num{m}, (G \, \num{n} \,\num{m} \, \num{h(n,m)})}\\
    & \red \tuple{\num{m+1}, \num{g(n, m, h(n, m))}}
  \end{align*}
  Como $g(n, m, h(n, m)) = h(n, m+1)$, concluímos a prova.

  Por fim, considere
  \begin{align*}
    H\,\num n\,\num m &\ident \lambd[x][\lambd[y][\fn{Snd} (y
        (\lambd[p].\tuple{\fn{Succ} (\fn{Fst}\, p), (G\, x\,
          (\fn{Fst}\, p)\, (\fn{Snd}\, p))}) \tuple{\num{0}, F x})]]\\
    & \qquad\,\num n\, \num m\\
    & \red \fn{Snd} (\num m \,
    \underbrace{(\lambd[p].\tuple{\fn{Succ} (\fn{Fst}\, p), (G \,\num n\,
      (\fn{Fst}\, p) (\fn{Snd}\, p))})}_{D_n} \tuple{\num{0}, F \num n})\\
    & \ident \fn{Snd} (\num m \, D_n\, \tuple{\num{0}, F \num n})\\
    & \red \fn{Snd} \,(D_n^m  \tuple{\num{0}, F \num n})\\
    & \red \fn{Snd} \,\tuple{\num{m}, \num{h(n, m)}} \\
    & \red \num{h(n,m)}. 
  \end{align*}  
\end{proof}


\begin{prop}
  Toda função recursiva primitiva é !!{lambda definable}.
\end{prop}

\begin{proof}
  Por \olref{lem:basic}, todas as funções básicas são !!{lambda definable},
  e, por \olref{lem:comp} e \olref{lem:prim}, as funções !!{lambda definable}
  são fechadas sob composição e recursão primitiva.
\end{proof}

\end{document}
